% !TEX root = ../journal_0421052099.tex
%
% Condensed from the error-analysis half of thesis_book/chapters/evaluation.tex.
% The histogram is the book's tables/tab_failure_histogram.tex, which
% scripts/build_figures.py generates from results/phase4/thesis_numbers.json.
\section{Error Analysis}\label{sec:errors}

\Cref{sec:results} measured how often each system fails. This section reads the
failures and names the mechanism of each one. The census is a population, not a
sample. Every question the full system did not answer correctly was read and
labelled by hand, after the confirmed benchmark defects were removed, for $256$
failures in all. The reading surfaced $16$ more benchmark defects inside the
census. They are filed as gold noise or ambiguous questions, so those two rows
of \Cref{tab:census} count the benchmark's failures rather than the system's.
Each failure packet holds the question, the gold answer, the answer given, the
plan, the explored relations with their scores, the backtrack reasons, and the
verifier outcome. The protocol was fixed before reading began. For each packet,
find where the trajectory first left the correct path, and label that point
rather than the last visible symptom. Without that rule, a run that
mis-decomposes, then explores the wrong relation, then hedges would be counted
three times.

\subsection{The Shape of the Failures}

\begin{table}[!t]
\caption{The failure census of AGR, $n = 256$. Each failure is labelled at the
  point where its trajectory first left the correct path. Wrong answers and
  hedges are kept separate.}
\label{tab:census}
\centering
\small
\setlength{\tabcolsep}{3pt}
\begin{tabular}{@{}lrrrrr@{}}
\toprule
& \multicolumn{2}{c}{WebQSP} & \multicolumn{2}{c}{CWQ} & \\
\cmidrule(lr){2-3} \cmidrule(lr){4-5}
Category & Wrong & Hedge & Wrong & Hedge & Total \\
\midrule
Relation selection & $19$ & $6$ & $20$ & $19$ & $64$ \\
Composite claim & $1$ & --- & $23$ & $23$ & $47$ \\
Knowledge-graph gap & $8$ & $4$ & $7$ & $24$ & $43$ \\
Decomposition error & $15$ & $11$ & $10$ & $2$ & $38$ \\
Answer selection & $1$ & $1$ & $8$ & $5$ & $15$ \\
Echo attractor & $5$ & $2$ & $4$ & $2$ & $13$ \\
Ambiguous question & $1$ & $1$ & $3$ & $5$ & $10$ \\
Premature termination & $1$ & $4$ & $1$ & $2$ & $8$ \\
Gold noise & $1$ & --- & $2$ & $3$ & $6$ \\
Other & --- & $1$ & --- & $5$ & $6$ \\
Verifier, wrongly rejected & --- & $3$ & $1$ & $1$ & $5$ \\
Verifier, wrongly accepted & --- & --- & --- & $1$ & $1$ \\
\midrule
Total & $52$ & $33$ & $79$ & $92$ & $256$ \\
\bottomrule
\end{tabular}
\end{table}

\Cref{tab:census} gives the histogram. Relation selection sits at or near the
top on both datasets, with $25$ of WebQSP's $85$ failures and $39$ of CWQ's
$171$. Picking the wrong edge is the base rate of graph navigation. What
distinguishes the datasets is a pair of categories that are nearly absent from
one and dominant in the other. Composite claims account for $1$ failure on
WebQSP and $46$ on CWQ, and knowledge-graph gaps for $12$ and $31$.
Decomposition errors run the other way, at $26$ against $12$.

Both asymmetries follow from how the datasets were built. WebQSP is $64\%$
one-hop, so it presents the planner with exactly the questions that do not need
decomposing. That mismatch is the significant ablation result of
\Cref{sec:ablations}, appearing here as a category count. CWQ's templates invoke
dates, numbers, and ordinals constantly, and the environment can express none of
them (\Cref{sec:gate}). When those gaps make a CWQ question unanswerable, AGR
more often declines than errs, with $24$ hedges against $7$ wrong answers in
that category.

\subsection{Failures inside the Agent}

\textbf{The entity-extraction bug.} Nine decomposition errors share one
mechanism, and they are the most actionable finding in the census. Asked ``who
was stephen r covey'', the system drafts ``Stephen Covey was a Writer,
Motivational speaker, Author, Consultant, Professor, and Manager''. All six gold
values are there. But the scored entity list comes back as the sentence's
grammatical subject, \textsf{[Stephen Covey]}. Navigation and verification both
succeeded. The step that turns answer text into an entity list failed. So the
measured accuracy on this question shape is a floor.

\textbf{Right candidate, wrong draft.} In $15$ answer-selection failures the
exact gold value entered the resolved candidates, and the drafter chose
something else. One was first filed as a verifier error. The draft read ``Marlon
Brando also played in \emph{Joy}'', and the verifier rejected it. Brando was not
in \emph{Joy}, but De Niro was, and the evaluator had resolved him. The verifier
was correct, and the case is the clearest example in the census of the layer
doing its job. A related family of six shares a mechanism nothing predicted. The
evaluator resolves the gold value, the verifier returns \textsf{grounded}, and
the draft hedges anyway. One answer states that Spanish was spoken in Spain and
then that the first language spoken there could not be determined.

\textbf{Context-stripping.} Three decomposition failures show a mechanism that
decomposition's own design cannot see. Splitting a question into steps removes a
qualifier that appears only in a later clause. So an early sub-objective
resolves the wrong sense of an ambiguous term before the disambiguating
information arrives. The strongest specimen asks for the 33rd president who
``was the nation's leader during WW2''. That is one entity under two
descriptions. The planner emits two unlinked sub-objectives, as though the
descriptions were a chain. With the first one active, the anchor \textsf{World
War II} scores at the noise floor, and the war constraint is never used. The run
drafts Woodrow Wilson, the 28th president. The undecomposed run answers Harry
S.~Truman correctly in three calls.

\textbf{Over-decomposition.} The more expected decomposition failures appear
too. In one, the plan splits ``what type of religion did massachusetts have''
into a step that finds Massachusetts and a step that finds its religion type.
The run burns $14$ of its $16$ calls on three evaluator backtracks and hedges,
while the undecomposed run finds the same fact in three calls. In another, the
planner rewrites ``where is ancient phoenician'' as the single step ``find
ancient Phoenician''. The rewrite drops the interrogative, the word that says
what is being asked. So the explorer re-identifies the topic instead of locating
it, and the run answers ``Phoenicia'', grounded and useless. With the ask
removed, the nearest well-grounded entity to a topic is the topic itself.

\textbf{Systemic relation gaps.} Relation selection is the largest category, at
$64$ failures. Either a wrong relation was explored, or the right one was
explored and abandoned. Its interest lies in cases that show the same gap from
several angles. Three questions on CWQ share a stem and reach Beyonc\'{e} by
three different routes, through a documentary, an actor, and a composer. All
three resolve her correctly. None ever tries \textsf{people.person.children},
the single relation that answers each of them, and all three hedge. Three
independent entry points hitting the same dead end show that the gap is
systemic. It is a property of the scorer's ranking over that entity's relations,
not of one question's phrasing.

\textbf{Abandoned relations.} Five premature terminations share one pattern
across unrelated domains. The correct relation is explored with a solid score,
as high as $0.731$, and the evaluator backtracks away from it and never returns.
That reads as a policy defect, not a reasoning gap.

\textbf{Verifier errors.} The two error polarities are very unevenly evidenced,
and the imbalance comes from the instrumentation. Five failures came from a
wrongly rejected claim, all by one mechanism. Each claim was true, but no single
edge stated it in the phrasing the drafter used, and the retry lost the answer.
Three of the five had already resolved the exact gold answer before the
rejection forced the retry. The verifier ran $828$ times over the $800$ test
questions, since a repaired question is verified again. Its first verdict was
unsupported on $52$ questions. After repairs, $761$ questions were answered
under a grounded verdict and $39$ under an unsupported one. So roughly one
rejected question in ten cost a correct answer, and even that is a lower bound.
Wrongful acceptance has no logged population at all, because accepted claims are
not persisted. One specimen surfaced by hand. A claim connecting Seth MacFarlane
to \emph{ThunderCats} was accepted, and a direct Cypher query confirms that no
such edge exists. A second case shows why acceptance is hard to judge. Asked
when President Wilson was in office, the system answered with inauguration
events whose edges are real, since the environment holds no dates. The
structural check passed correctly while the answer is wrong. A check of whether
a claim is \emph{in} the graph cannot check whether it is \emph{the answer}.

\subsection{Knowledge-Graph Gaps}

Knowledge-graph gaps are the third-largest category, at $43$ failures, and CWQ's
largest hedge category. They are the unanswerable class of \Cref{sec:gate} seen
in the census, and five subtypes account for $39$ of them. Ordinal questions
($9$) ask for a first, last, or smallest member. The relevant relation is
explored and the candidates come back, but no component compares them. Asked for
the year of the first Miss America pageant, the system answers $1925$, then
$1980$, then $1929$ across successive rounds, because nothing sorts and takes
the extreme. Numeric literals ($16$) are internal identifiers and codes that
CWQ's templates invoke and that no explored relation exposes. Date literals
($4$) and temporal qualifiers ($4$) fail because the environment holds events
but not their dates. Asked who was president in 1988, the system names the
winner of that November's election, who took office the following January. Only
the last subtype, data errors ($6$), is a gap in this particular import. The
other four are gaps in what a triple-and-traversal architecture can express at
all. No larger graph fixes them, since they need an operator.

\subsection{The Echo Attractor}\label{sec:echo}

Thirteen failures name a mechanism that recurs independently of any component.
The system resolves a real, well-grounded entity that sits near the correct
answer in the graph. The verifier passes it, because it genuinely is grounded.
The fact is true but not what was asked. Two shapes appear. In the first, the
system confuses entities in one tight cluster. Asked who plays Lex Luthor on
\emph{Smallville}, it answers that John Glover plays Lionel Luthor. In the
second, it finds the right entity family at the wrong level. It answers
``Kingdom of Denmark'' where the gold answer is ``Denmark'', or narrows from
Harvard College to Harvard University. None of these are cases where the system
knew nothing. The right edge exists, is used, and lands one node away.

The echo attractor is also why the consensus pass of \Cref{sec:gold} flagged far
more questions than adjudication confirmed. All five systems answer ``rick
scott'' where the gold answer lists his professions, and ``thor heyerdahl''
where it lists his occupations. Each returns the topic entity. All five answer
``belgium'' where the gold answer is the containing region, and ``maryland''
where it lists counties. When systems agree with one another against gold, the
usual cause is a shared attractor, not a label error. A policy of rescoring
whenever three systems agree would have turned this failure into apparent
correctness. Adjudication is what stands between the two. Counting the census
and the earlier exclusions together, $57$ distinct questions needed the
benchmark, not the system, to be corrected. Two patterns account for most of the
CWQ defects, and both are properties of template generation rather than of
annotation care. One is a gold list flattened from a relation that has several
values, with the pipeline keeping a marginal one. The other is an annotation
query that traversed one relation too far.

\endinput
